\section{Fórmula de inversión}

Sea $f\in\mathcal S(\mathbb R^2)$. Partimos de la fórmula de inversión de Fourier en dos variables:
\[
f(x)=
\int_{\mathbb R^2}
\widehat f(\xi)e^{2\pi i x\cdot\xi}\,d\xi.
\]
Para reescribir esta integral en coordenadas adaptadas a la transformada de Radon, usamos el cambio con radio firmado
\[
\Phi(\sigma,\theta)=\sigma\omega(\theta),
\qquad
\sigma\in\mathbb R,
\quad
\theta\in[0,\pi).
\]
Para cada $\xi\neq0$ existe una única pareja
$(\sigma,\theta)\in(\mathbb R\setminus\{0\})\times[0,\pi)$
tal que $\xi=\sigma\omega(\theta)$. El origen es el único punto
degenerado de la parametrización y no afecta la integral.

Las derivadas del cambio son
\[
\partial_\sigma\Phi=\omega(\theta),
\qquad
\partial_\theta\Phi=\sigma\omega^\perp(\theta).
\]
Por tanto,
\[
\left|
\det\left(
\partial_\sigma\Phi,
\partial_\theta\Phi
\right)
\right|
=
|\sigma|.
\]
Así,
\[
d\xi=|\sigma|\,d\sigma\,d\theta.
\]

Sustituyendo en la inversión de Fourier y usando el teorema de corte,
\[
\begin{aligned}
f(x)
&=
\int_0^\pi\int_{\mathbb R}
\widehat f(\sigma\omega(\theta))
e^{2\pi i\sigma x\cdot\omega(\theta)}
|\sigma|\,d\sigma\,d\theta \\
&=
\int_0^\pi\int_{\mathbb R}
\mathcal F_t(\mathcal Rf)(\sigma,\theta)
|\sigma|
e^{2\pi i\sigma x\cdot\omega(\theta)}
\,d\sigma\,d\theta.
\end{aligned}
\]
Esta es la fórmula de inversión que se usará en la tesis. Con la convención de Fourier fijada y con integración angular en $[0,\pi)$, no aparece un factor adicional $1/2$ ni $1/\pi$ en esta expresión. Otras fuentes pueden usar frecuencia angular o integración sobre todas las direcciones, lo que modifica las constantes. La inversión de la transformada de Radon se desarrolla en Kuchment \cite[Sec.~5.7, pp.~36--39]{kuchment2014} y en Natterer \cite[Sec.~II.2, Thm.~2.1, pp.~18--20]{natterer2001}, con convenciones que deben convertirse antes de comparar factores.
